\section{Course Philosophy and the Three-Document Ecosystem}

Capstone Design is not another course; it is a professional practice experience compressed into two semesters. You will define problems, manage a team, communicate with a real client, design and build (or model) a system, verify it against requirements, and present your results publicly at the Design Showcase. The skills you develop here (engineering judgment, project management, client communication, teamwork) are the skills employers value most.

\begin{faqbox}[What level of independence is expected?]
Capstone demands self-motivation. Unlike lecture-based courses, your progress depends on work you organize and complete \emph{outside} of scheduled Tuesday/Thursday class sessions. Your PA will not chase you for deliverables. Your client will not remind you of deadlines. You and your teammates are responsible for holding each other accountable, scheduling your own work sessions, and maintaining momentum between class meetings. Treat this course the way you would treat your first engineering job: your supervisor sets the objectives, but you manage your own time to meet them.
\end{faqbox}

\subsection{How SDI and SDII Fit Together}
\textbf{SDI} (EDNS~491) is the \textbf{definition semester}: you go from project assignment through requirements, concept selection, and Preliminary Design Review. Sprints~0--6 cover the left side of the V-Model.

\textbf{SDII} (EDNS~492) is the \textbf{execution semester}: you re-engage with the client, complete the Critical Design Review, build/model your system, test it, and present at the Final Design Review and Design Showcase. Sprints~7--14 cover the bottom and right side of the V-Model.

\subsection{Your Three Documents}
\begin{itemize}[nosep,leftmargin=1.5em]
\item \textbf{This Student Guide}: tells you what to focus on each sprint.
\item \textbf{The Course Text}: provides the engineering methodology, process guidance, and technical reference. Organized by topic in 22 chapters across 5 parts.
\item \textbf{Canvas}: assignment rubrics, submission portals, professional development modules, scheduling, and the Labriola InnoHub enrollment site.
\end{itemize}


% ═══════════════════════════════════════
